\bookchapter{Bab 33: Prospek Masa Depan — Perjalanan Berikutnya Ponsel}

\chapterintro{Ponsel tidak akan hilang semalam, tapi “dominasi absolutnya” mungkin dialihkan oleh perangkat generasi baru. Bab ini kita bahas tiga hal: seperti apa perangkat generasi berikutnya, bagaimana ponsel “menyembunyikan diri” menjadi fondasi identitas baru, dan isu etika yang harus kita pikirkan lebih dulu.}

\booksection{AR/VR dan Spatial Computing: “Desktop Kedua” di Luar Ponsel}

Persaingan perangkat generasi berikutnya sudah dimulai, yang paling mungkin mengalihkan waktu ponsel adalah perangkat AR/VR dan spatial computing:

\begin{itemize}
  \item \textbf{Vision Pro generasi berikutnya}: Rilis Apple Vision Pro menandai pendaratan resmi konsep "spatial computing", tapi harga produk generasi pertama (29999 yuan+) dan berat masih menjadi ambang. Diperkirakan 2026-2027 akan merilis versi Vision yang lebih ringan dan murah, membuat spatial computing dari "mainan developer" menuju "early majority". Saat itu iPhone mungkin menjadi pendamping komputasi dan pusat data bagi kacamata Vision.
  \item \textbf{Kacamata AR}: Buka "desktop kedua" di depan mata Anda, informasi digital menjadi seperti udara ada di mana-mana. Kerjasama Meta dan Ray-Ban smart glasses sudah jual jutaan unit, dalam negeri Rokid, XREAL, INMO dan pabrikan lain juga merilis AR glasses konsumen. AR glasses 2025-2026 sudah bisa real-time translation, HUD navigasi, AI assistant conversation, tapi masih terbatas oleh tampilan waveguide dan battery life, jarak "menggantikan ponsel" masih 3-5 tahun.
  \item \textbf{Smart ring dan wearable}: Samsung Galaxy Ring, Oura Ring dan smart ring lain mulai populer, sebagai "sensor tempelan" ponsel melengkapi monitoring kesehatan. Peran ponsel sedang bergeser dari "terminal serbaguna" ke kombinasi "hub + multiple aksesori ringan".
\end{itemize}

\booksection{Brain-Computer Interface: Pikiran Menjadi Media Interaksi}

Pikiran itu sendiri akan menjadi media interaksi, tapi jalan ini lebih panjang dari yang dibayangkan:

\begin{itemize}
  \item \textbf{Perkembangan terbaru Neuralink}: Neuralink milik Musk telah menyelesaikan implantasi manusia pertama pada 2024, subjek pertama Noland Arbaugh mengontrol mouse komputer dengan pikiran untuk bermain catur. Saat ini BCI (brain-computer interface) implan terutama untuk skenario bantuan pasien lumpuh, tujuannya membantu mereka kembali "mengoperasikan" dunia digital.
  \item \textbf{BCI non-invasif}: Beberapa perusahaan mengembangkan headband non-invasif (seperti Kernel, OpenBCI), tanpa operasi bisa membaca sinyal otak, tapi presisi jauh di bawah invasif. Masih cukup jauh dari "mengetik dengan pikiran".
  \item \textbf{Prospek konsumen}: Brain-computer interface masih 10-20 tahun dari interaksi ponsel konsumen, tapi jangka pendek mungkin lebih dulu dipakai untuk skenario bantu aksesibilitas, monitoring fokus dll. Isu etika dan privasi ("Gelombang otak saya termasuk privasi siapa") akan menjadi fokus perdebatan sebelum teknologi matang.
\end{itemize}

\booksection{“Menyembunyikan Diri” Ponsel dan Fondasi Identitas Baru}

Arah evolusi masa depan ponsel yang lain mungkin adalah “menyembunyikan diri” — berubah menjadi “dalang di balik layar” yang ada di mana-mana, menjadi satu-satunya fondasi identitas digital Anda:

\begin{itemize}
  \item \textbf{Autentikasi dua faktor berbasis cloud}: Seketika selesaikan verifikasi "Anda adalah Anda", login layanan apa pun tanpa input manual kode verifikasi.
  \item \textbf{Warisan digital menjadi pengaturan pabrik}: Saat aktivasi ponsel baru preset skema pewarisan aset digital, membuat "urusan setelah meninggal" menjadi default pabrik.
  \item \textbf{Dari "pusat perangkat" ke "pusat identitas"}: Brankas data pribadi cloud Anda adalah pusatnya, hardware hanyalah akses masuk. Ganti ponsel, ganti perangkat, identitas dan data selalu mengikuti Anda.
\end{itemize}

\booksection{Refleksi Etika: Apakah Kita Masih Tuan dari Ponsel?}

Satu pertanyaan ultimate selalu ada: Apakah kita “dewa” yang diberdayakan teknologi, atau “otak dalam tank” yang diberi makan algoritma?

\begin{itemize}
  \item \textbf{Etika AI}: Saat on-device AI lebih paham preferensi Anda dari diri Anda sendiri, "pilihan" Anda adalah kehendak mandiri atau hasil yang dipandu? Saat AI Agent memesan dan membalas pesan untuk Anda, di mana tanggung jawab dan batasannya?
  \item \textbf{Kedaulatan data}: Data kesehatan, data gelombang otak, data biometrik Anda, milik Anda, milik pabrikan perangkat atau milik penyedia cloud? "Personal Information Protection Law" dan EU GDPR hanya awal, "hak portabilitas data" dan "hak untuk dilupakan" implementasi nyatanya masih jauh.
  \item \textbf{Kesenjangan digital}: Ponsel AI, folding screen, satelit communication semua memperlebar kesenjangan antara "mampu beli" dan "tidak mampu beli". Akses yang adil ke infrastruktur digital akan menjadi isu sosial baru — sama kritisnya dengan "bisa internet atau tidak" 20 tahun lalu.
  \item Saat algoritma lebih paham apa yang ingin Anda lihat daripada diri Anda sendiri, seberapa banyak "suka" Anda yang spontan?
  \item Saat Anda menyalakan layar tanpa sadar setiap beberapa menit, apakah ini kebiasaan Anda atau "perilaku adiktif" yang dirancang?
\end{itemize}

\booksection{Tetap Terhubung, Juga Tetap Mandiri — Filosofi Kehidupan Digital Anda}

Apa pun teknologi berevolusi, beberapa filosofi kehidupan digital berikut, 10 tahun lagi tetap berlaku:

\begin{enumerate}
  \item Jadilah pengalokasi sumber daya, bukan penerima push.
  \item Buat otomatisasi menjadi "co-pilot" setia Anda.
  \item Bangun "garis belakang digital" yang kokoh.
  \item Secara berkala lakukan "digital decluttering" dan "puasa sosial".
  \item Selalu jadikan keamanan dan privasi sebagai prasyarat semua kemudahan.
\end{enumerate}

\begin{tcbpractice}{Mari Coba: Susun "Syarat dan Ketentuan Layanan Penggunaan Perangkat Digital Pribadi" Anda}
Tulis satu “service agreement” sederhana untuk Anda dan ponsel Anda (serta AR glasses, AI assistant masa depan), untuk mendefinisikan hubungan kalian. Misalnya:

\begin{itemize}
  \item Pasal 1: Waktu perhatian saya adalah milik saya. Anda harus secara aktif masuk ke mode "Jangan Ganggu" saat saya bekerja/bersama keluarga.
  \item Pasal 2: Data pribadi saya adalah properti saya. Tanggung jawab utama Anda adalah enkripsi yang aman dan sinkronisasi yang nyaman.
  \item Pasal 3: Tujuan bersama kita adalah meningkatkan efisiensi dan pengalaman hidup. Jika Anda membuat saya kecanduan melalui desain yang manipulatif, saya berhak meng-uninstall Anda.
\end{itemize}

\end{tcbpractice}

\begin{tcbsummary}{Ringkasan Bab}
\begin{itemize}
  \item Vision Pro dan AR glasses adalah new computing entry point di luar ponsel, tapi sulit menggantikan ponsel jangka pendek
  \item Brain-computer interface masih 10-20 tahun dari konsumen, isu etika sebelum teknologi matang
  \item Ponsel sedang berevolusi ke arah "menyembunyikan diri", dari device-centric menuju identity-centric
  \item Etika AI, kedaulatan data, kesenjangan digital adalah isu inti 10 tahun ke depan
  \item Tetap terhubung juga tetap mandiri, adalah filosofi kehidupan digital yang tidak berubah
\end{itemize}

\end{tcbsummary}
